\documentclass[11pt,a4paper]{article}

% ============================= PAQUETES ===================================
\usepackage[utf8]{inputenc}
\usepackage[T1]{fontenc}
\usepackage[spanish,es-tabla]{babel}
\usepackage[margin=2.5cm]{geometry}
\usepackage{amsmath}
\usepackage{graphicx}
\usepackage{booktabs}
\usepackage{longtable}
\usepackage{array}
\usepackage{caption}
\usepackage{float}
\usepackage{xcolor}
\usepackage{hyperref}
\usepackage{pdflscape}
\usepackage{enumitem}
\usepackage{fancyhdr}
\usepackage{titlesec}
\usepackage[numbers]{natbib}

\hypersetup{
    colorlinks=true,
    linkcolor=blue!50!black,
    citecolor=blue!50!black,
    urlcolor=blue!50!black
}

\titleformat{\section}{\normalfont\Large\bfseries}{\thesection.}{0.6em}{}
\titleformat{\subsection}{\normalfont\large\bfseries}{\thesubsection.}{0.6em}{}

\pagestyle{fancy}
\fancyhf{}
\setlength{\headheight}{14.5pt}
\rhead{Informe Final -- Análisis Exploratorio de Datos}
\lhead{UNSA -- Tópicos en Ciencia de Datos}
\rfoot{\thepage}

\newcolumntype{P}[1]{>{\raggedright\arraybackslash}p{#1}}

% ============================= PORTADA =====================================
\begin{document}

\begin{titlepage}
\centering
\vspace*{1.5cm}

{\large Universidad Nacional de San Agustín de Arequipa\par}
{\large Maestría en Ciencia de la Computación\par}
{\large Curso: Tópicos en Ciencia de Datos\par}
\vspace{0.3cm}
{\small Docente: Ana María Cuadros Valdivia\par}

\vspace{2.5cm}
{\Huge\bfseries Informe Final\par}
\vspace{0.5cm}
{\Large Análisis Exploratorio de Datos\par}
\vspace{0.8cm}
{\large\itshape Patrones espacio-temporales de criminalidad en Chicago (2017)\par}
{\large\itshape como base para una propuesta de Visual Analytics\par}

\vspace{2.5cm}
{\large\bfseries Estudiante:\par}
{\large Alex Maglio Neyra Herrada\par}

\vspace{2cm}
{\large Fecha de presentación: 19 de septiembre de 2026\par}

\vfill
{\small Paper base del proyecto final de Visual Analytics:\par
Garcia-Zanabria et al. (2021), \textit{CrimAnalyzer: Understanding Crime
Patterns in São Paulo}, IEEE TVCG.\par}
\end{titlepage}

\tableofcontents
\newpage

% ============================================================================
\section{Definición del problema}

El problema que se aborda es la dificultad que enfrentan los analistas de
seguridad pública y los tomadores de decisión de una ciudad para identificar,
de forma rápida e interactiva, \textbf{dónde y cuándo se concentran los
distintos tipos de delito}, y qué factores contextuales (tipo de lugar, hora
del día, día de la semana) se asocian a una mayor o menor probabilidad de que
un incidente termine en arresto. Tradicionalmente, este análisis se apoya en
mapas de calor estáticos que muestran la intensidad del crimen en un momento
dado, pero que no permiten explorar cómo cambian esos patrones en el tiempo
ni comparar el comportamiento de distintos subtipos de delito de forma
interactiva.

Resolver este problema es relevante en varios niveles:

\begin{itemize}
    \item \textbf{Social:} una mejor comprensión de los patrones delictivos
    contribuye a políticas de seguridad ciudadana más focalizadas.
    \item \textbf{De gestión pública / negocio:} los departamentos de policía
    operan con recursos de patrullaje limitados; identificar hotspots
    espaciales y temporales permite asignarlos de manera más eficiente.
    \item \textbf{Científica:} el trabajo se enmarca en la línea de
    investigación de \textit{Visual Analytics} aplicada a criminalidad, cuyo
    referente directo es \textbf{CrimAnalyzer} \cite{crimanalyzer}, un sistema
    visual-interactivo desarrollado para estudiar hotspots de crimen y su
    evolución temporal en São Paulo mediante factorización de matrices no
    negativas (NMF) en lugar de los métodos estadísticos tradicionales de
    mapeo de calor (KDE, Getis-Ord Gi*).
\end{itemize}

El área de conocimiento involucrada es \textbf{Visual Analytics aplicado a
Seguridad Pública / Crime Analytics}, dentro del campo más amplio de la
Ciencia de Datos.

Al finalizar el proyecto se espera contar con un sistema de Visual Analytics
--inspirado en la metodología de CrimAnalyzer pero aplicado a un dataset
público de la ciudad de Chicago-- que permita a un analista explorar de
forma interactiva la concentración espacial del crimen (hotspots), su
evolución temporal y su relación con variables como el tipo de lugar y la
tasa de esclarecimiento (arresto). El presente informe constituye la etapa
de \textbf{análisis exploratorio} que sienta las bases de conocimiento y de
calidad de datos necesarias para diseñar dicho sistema.

% ============================================================================
\section{Hipótesis / Preguntas de investigación}

\subsection{Motivación}

Las hipótesis de este informe surgen de dos fuentes complementarias:

\begin{enumerate}
    \item El artículo científico base, \textbf{CrimAnalyzer}
    \cite{crimanalyzer}, que identifica tres desafíos centrales al analizar
    crimen en una ciudad: (i) la necesidad de explorar patrones
    \emph{espaciales} locales (hotspots), (ii) identificar hotspots que no
    necesariamente son los de mayor volumen pero sí relevantes para la
    investigación criminal, y (iii) comprender la dinámica \emph{temporal}
    de esos patrones.
    \item Conocimiento de dominio en criminología (estacionalidad horaria de
    los delitos, diferencias en la tasa de esclarecimiento según el tipo de
    delito, y la asociación conocida entre violencia doméstica y el hogar
    como lugar de ocurrencia).
\end{enumerate}

A partir de esto se formulan 4 hipótesis (se exige un mínimo de 3), cada una
directamente respondible con las variables disponibles en el dataset de
Chicago: \texttt{Primary Type}, \texttt{Date}, \texttt{District},
\texttt{Location Description}, \texttt{Arrest} y \texttt{Domestic}.

\subsection{Hipótesis formuladas}

\begin{description}[leftmargin=1.6cm, style=nextline]
    \item[Hipótesis 1 (espacial):] ¿Se concentra el crimen de forma desigual
    entre los distritos policiales de Chicago, evidenciando un patrón de
    \emph{hotspots} espaciales comparable al reportado por CrimAnalyzer para
    São Paulo?

    \item[Hipótesis 2 (temporal):] ¿El volumen de delitos varía de forma
    significativa según la hora del día y el día de la semana, mostrando
    patrones identificables de picos y valles horarios?

    \item[Hipótesis 3 (esclarecimiento):] ¿La probabilidad de que un
    incidente termine en arresto depende del tipo de delito
    (\texttt{Primary Type})?

    \item[Hipótesis 4 (violencia doméstica y lugar):] ¿La proporción de
    incidentes clasificados como violencia doméstica varía significativamente
    según el tipo de lugar donde ocurre el delito
    (\texttt{Location Description})?
\end{description}

% ============================================================================
\section{Plan de análisis}

El análisis sigue las etapas estándar del ciclo de vida de Ciencia de Datos
aplicadas al pipeline de AED:

\begin{enumerate}
    \item \textbf{Obtención de datos:} descarga del extracto oficial del
    dataset "Crimes -- 2001 to Present" de la Ciudad de Chicago
    correspondiente al período enero--agosto de 2017.
    \item \textbf{Wrangling:} parseo de fechas, creación de variables
    derivadas (\texttt{Hour}, \texttt{DayOfWeek}, \texttt{Month}),
    recodificación de variables booleanas.
    \item \textbf{Limpieza (cleaning):} cuantificación y tratamiento de
    valores nulos, detección de valores fuera de dominio, verificación de
    duplicados.
    \item \textbf{Exploración (AED):} estadística descriptiva univariada y
    bivariada, análisis de correlación.
    \item \textbf{Visualización:} generación de gráficos (histogramas,
    boxplots, barras, heatmaps, violin plot, mapa de dispersión y dashboard).
    \item \textbf{Respuesta a hipótesis:} contraste de cada hipótesis con
    evidencia visual y pruebas estadísticas ($\chi^2$ de independencia).
    \item \textbf{Conclusiones y puente a Visual Analytics:} identificación
    de los hallazgos y problemas de calidad prioritarios para el diseño del
    sistema interactivo.
\end{enumerate}

% ============================================================================
\section{Fuente de datos}

\textbf{Origen.} El dataset se obtiene del portal de datos abiertos de la
Ciudad de Chicago, \textit{Chicago Data Portal}, en el conjunto
\href{https://data.cityofchicago.org/Public-Safety/Crimes-2001-to-Present/ijzp-q8t2}{\texttt{Crimes -- 2001 to Present}}
(identificador \texttt{ijzp-q8t2}) \cite{chicagoportal}. Es publicado y
mantenido por el \textbf{Departamento de Policía de Chicago (CPD)} y se
extrae directamente de su sistema operacional \textbf{CLEAR} (\textit{Citizen
Law Enforcement Analysis and Reporting}). El portal se actualiza a diario y
excluye los últimos 7 días para permitir la verificación de los reportes.
Para este informe se utilizó el extracto correspondiente al período
\textbf{1 de enero al 25 de agosto de 2017} (172,327 registros), obtenido
como réplica exacta del esquema oficial desde un repositorio público de
datos abiertos en GitHub que re-expone el dataset de la ciudad
\cite{githubmirror}.

\textbf{Área de conocimiento.} Seguridad pública / criminología
computacional, con técnicas de Ciencia de Datos y Visual Analytics.

\textbf{Problema que se quiere resolver.} A nivel de negocio: apoyar la
asignación de recursos de patrullaje y la comprensión de patrones
delictivos por parte de analistas y tomadores de decisión. A nivel
computacional: caracterizar la distribución espacio-temporal del crimen y
sus factores asociados (lugar, tipo, esclarecimiento) a partir de datos
tabulares de incidentes.

\textbf{Por qué se capturaron las variables.} Los campos del dataset
corresponden a los datos que el sistema CLEAR registra por razones
administrativas y legales de cada denuncia: identificación del caso,
clasificación del delito (\texttt{IUCR}, \texttt{Primary Type},
\texttt{Description}, \texttt{FBI Code}), ubicación aproximada
(\texttt{Block}, \texttt{Beat}, \texttt{District}, \texttt{Ward},
\texttt{Community Area}, coordenadas), fecha/hora del incidente, y dos
banderas de interés investigativo (\texttt{Arrest}, \texttt{Domestic}). La
ubicación se publica únicamente a nivel de cuadra (\textit{block}) para
proteger la privacidad de las víctimas.

\textbf{Referencia del paper base.} El proyecto de Visual Analytics de este
curso toma como referencia metodológica \textbf{CrimAnalyzer}
\cite{crimanalyzer}, que originalmente aplicó un enfoque de detección de
hotspots basado en NMF sobre datos policiales reales de São Paulo. El
presente informe traslada esa misma pregunta de investigación --y una
metodología equivalente de análisis exploratorio-- a un dataset de acceso
público (Chicago), evitando así las restricciones de acceso del dataset
original de São Paulo, que no está disponible públicamente.

% ============================================================================
\section{Formato original del dataset}

El dataset se distribuye en formato \textbf{CSV} (también disponible en
JSON y GeoJSON a través de la API abierta Socrata/SODA del portal, y
mediante el paquete de Python \texttt{sodapy}). Particularidades encontradas
al leerlo:

\begin{itemize}
    \item El campo \texttt{Date} viene como \textbf{texto} con el formato
    \texttt{MM/DD/AAAA HH:MM:SS AM/PM}, por lo que requiere un parseo
    explícito (\texttt{pd.to\_datetime}) antes de poder derivar variables
    temporales; no llega tipado como fecha nativa.
    \item Los campos \texttt{Arrest} y \texttt{Domestic} sí llegan
    correctamente tipados como booleanos (\texttt{true}/\texttt{false}).
    \item Los registros sin geocodificación no traen un marcador explícito de
    nulo, sino que las columnas de coordenadas (\texttt{X Coordinate},
    \texttt{Y Coordinate}, \texttt{Latitude}, \texttt{Longitude},
    \texttt{Location}) aparecen vacías, y \texttt{pandas} las infiere
    correctamente como \texttt{NaN} al cargarlas.
    \item Variables categóricas codificadas numéricamente
    (\texttt{District}, \texttt{Ward}, \texttt{Beat},
    \texttt{Community Area}) se cargan por defecto como tipo entero, lo cual
    puede inducir a tratarlas erróneamente como variables continuas si no se
    documenta su verdadera naturaleza nominal.
\end{itemize}

% ============================================================================
\section{Descripción del conjunto de datos}

\subsection{A nivel de atributos}

El dataset cuenta con \textbf{172,327 registros} y \textbf{22 atributos}. La
Tabla \ref{tab:atributos} resume cada uno de ellos.

\begin{landscape}
\begin{longtable}{P{2.6cm} P{4.0cm} P{1.6cm} P{2.0cm} P{4.4cm} P{3.4cm}}
\caption{Descripción de atributos del dataset \textit{Crimes -- 2017 (Chicago)}}
\label{tab:atributos}\\
\toprule
\textbf{Atributo} & \textbf{Significado} & \textbf{Tipo} & \textbf{Nulos (n / \%)} & \textbf{Rango / valores únicos} & \textbf{Relevancia}\\
\midrule
\endfirsthead
\toprule
\textbf{Atributo} & \textbf{Significado} & \textbf{Tipo} & \textbf{Nulos (n / \%)} & \textbf{Rango / valores únicos} & \textbf{Relevancia}\\
\midrule
\endhead
\bottomrule
\endfoot
ID & Identificador único del registro & Numérico discreto (id) & 0 / 0.00\% & 172,327 valores únicos & Baja (solo identificador) \\
Case Number & Número de caso asignado por CPD & Categórico nominal (texto) & 0 / 0.00\% & 172,294 únicos (27 casos de homicidio con 2--4 víctimas comparten número) & Media (granularidad de registros) \\
Date & Fecha y hora del incidente & Fecha/hora (texto a parsear) & 0 / 0.00\% & 01/01/2017 -- 25/08/2017 & Alta (H2) \\
Block & Dirección aproximada a nivel de cuadra & Categórico nominal & 0 / 0.00\% & 25,385 valores únicos & Baja (alta cardinalidad) \\
IUCR & Código de clasificación policial de Illinois & Categórico nominal & 0 / 0.00\% & 318 categorías & Media (subcategoría del delito) \\
Primary Type & Categoría general del delito & Categórico nominal & 0 / 0.00\% & 32 categorías (top: THEFT 40,742; BATTERY 32,394) & Alta (H1, H3) \\
Description & Subtipo textual del delito & Categórico nominal & 0 / 0.00\% & Cientos de valores (detalle de IUCR) & Media \\
Location Description & Tipo de lugar donde ocurrió & Categórico nominal & 515 / 0.30\% & 123 categorías (top: STREET 38,980) & Alta (H4) \\
Arrest & Si hubo arresto & Booleano & 0 / 0.00\% & \{True, False\}; tasa global = 19.24\% & Alta (H3) \\
Domestic & Si se clasificó como violencia doméstica & Booleano & 0 / 0.00\% & \{True, False\}; tasa global = 16.31\% & Alta (H4) \\
Beat & Subdivisión de patrullaje (anidada en District) & Categórico nominal (cod. numérico) & 0 / 0.00\% & 274 valores; rango 111--2535 & Media (redundante con District) \\
District & Distrito policial & Categórico nominal (cod. numérico) & 0 / 0.00\% & 23 valores (1--25; outlier administrativo 31 con 2 registros) & Alta (H1) \\
Ward & Distrito electoral / municipal & Categórico nominal (cod. numérico) & 1 / 0.00\% & 50 valores (1--50) & Media \\
Community Area & Área comunitaria oficial de Chicago & Categórico nominal (cod. numérico) & 0 explícitos & Rango 0--77 (válido 1--77; 1 registro con valor 0 fuera de dominio) & Media \\
FBI Code & Clasificación del delito según el FBI & Categórico nominal & 0 / 0.00\% & 26 categorías & Media \\
X Coordinate & Coordenada plana (State Plane, pies) & Numérico continuo & 6,102 / 3.54\% & media $\approx$ 1,164,934 pies & Media (soporte geoespacial) \\
Y Coordinate & Coordenada plana (State Plane, pies) & Numérico continuo & 6,102 / 3.54\% & -- & Media (soporte geoespacial) \\
Year & Año del registro & Numérico discreto & 0 / 0.00\% & Constante = 2017 & Nula (sin variabilidad) \\
Updated On & Fecha de última actualización del registro & Fecha/hora & 0 / 0.00\% & Cercana a ago--sep 2017 & Baja (control de versión, no analítica) \\
Latitude & Latitud (WGS84) & Numérico continuo & 6,102 / 3.54\% & 41.64° -- 42.02° & Alta (H1, mapa geográfico) \\
Longitude & Longitud (WGS84) & Numérico continuo & 6,102 / 3.54\% & $-$87.93° -- $-$87.52° & Alta (H1, mapa geográfico) \\
Location & Texto combinado "(lat, long)" & Categórico (texto redundante) & 6,102 / 3.54\% & Redundante con Latitude/Longitude & Nula (redundante) \\
\end{longtable}
\end{landscape}

\textbf{Medidas de posición y dispersión.} La Tabla \ref{tab:descriptivos}
presenta los estadísticos descriptivos de las variables numéricas /
codificadas más relevantes.

\begin{table}[H]
\centering
\caption{Estadísticos descriptivos de variables numéricas seleccionadas}
\label{tab:descriptivos}
\small
\begin{tabular}{lrrrrrrrr}
\toprule
Variable & Media & Desv. est. & Mín & Q1 & Mediana & Q3 & Máx & Asimetría\\
\midrule
Hour (hora, 0--23) & 13.21 & 6.53 & 0 & 9 & 14 & 19 & 23 & $-0.45$\\
Beat & 1142.51 & 698.24 & 111 & 611 & 1024 & 1713 & 2535 & 0.41\\
District & 11.20 & 6.97 & 1 & 6 & 10 & 17 & 31 & 0.41\\
Ward & 23.15 & 14.09 & 1 & 10 & 24 & 34 & 50 & 0.12\\
Community Area & 36.77 & 21.37 & 0 & 23 & 32 & 54 & 77 & 0.24\\
Latitude (°) & 41.84 & 0.09 & 41.64 & 41.77 & 41.86 & 41.91 & 42.02 & $-0.11$\\
Longitude (°) & $-87.67$ & 0.06 & $-87.93$ & $-87.71$ & $-87.66$ & $-87.63$ & $-87.52$ & $-0.34$\\
\bottomrule
\end{tabular}
\end{table}

Ninguna de estas variables numéricas presenta asimetría marcada (todos los
coeficientes están entre $-0.45$ y $0.41$); la variable \texttt{Hour}
muestra una leve asimetría negativa, coherente con la mayor concentración de
delitos en horas de la tarde/noche (ver Sección \ref{sec:hipotesis}). No se
identificaron outliers numéricos extremos en coordenadas geográficas: el
rango de latitud/longitud es consistente con los límites geográficos reales
de la ciudad de Chicago.

\subsection{A nivel de registros}

Cada fila del dataset representa \textbf{un incidente delictivo reportado} a
la policía de Chicago (es un \emph{reporte}, no necesariamente un delito
judicialmente confirmado). La variable \texttt{Primary Type} funciona como
la etiqueta o clase principal (32 categorías), mientras que \texttt{Arrest}
y \texttt{Domestic} son etiquetas binarias adicionales de interés
investigativo. La granularidad es, en general, de un registro por incidente,
\textbf{con una excepción documentada}: en los casos de \texttt{HOMICIDE} el
propio proveedor de datos indica que la granularidad es \emph{por víctima}
y no por caso; esto se confirmó empíricamente: 27 números de caso
(\texttt{Case Number}) están asociados a más de un registro (22 casos con 2
víctimas, 4 casos con 3 víctimas y 1 caso con 4 víctimas), y el 100\% de
esos 60 registros corresponden a \texttt{Primary Type = HOMICIDE}. La
ubicación se reporta a nivel de cuadra (\textit{block}), no de dirección
exacta, por motivos de privacidad de las víctimas.

\subsection{Relación entre atributos}

Se calcularon las correlaciones de \textbf{Pearson} y \textbf{Spearman}
entre las variables numéricas / codificadas disponibles (Tabla
\ref{tab:correlacion}); ambos coeficientes arrojaron valores muy similares.

\begin{table}[H]
\centering
\caption{Matriz de correlación de Pearson entre atributos numéricos}
\label{tab:correlacion}
\small
\begin{tabular}{lrrrrrrr}
\toprule
 & Beat & District & Ward & Com. Area & Hour & Arrest & Domestic\\
\midrule
Beat & 1.00 & 1.00 & 0.66 & $-0.49$ & 0.00 & $-0.04$ & $-0.06$\\
District & 1.00 & 1.00 & 0.66 & $-0.49$ & 0.00 & $-0.04$ & $-0.06$\\
Ward & 0.66 & 0.66 & 1.00 & $-0.53$ & 0.00 & $-0.03$ & $-0.08$\\
Community Area & $-0.49$ & $-0.49$ & $-0.53$ & 1.00 & 0.00 & 0.04 & 0.10\\
Hour & 0.00 & 0.00 & 0.00 & 0.00 & 1.00 & 0.05 & $-0.04$\\
Arrest & $-0.04$ & $-0.04$ & $-0.03$ & 0.04 & 0.05 & 1.00 & $-0.03$\\
Domestic & $-0.06$ & $-0.06$ & $-0.08$ & 0.10 & $-0.04$ & $-0.03$ & 1.00\\
\bottomrule
\end{tabular}
\end{table}

\textbf{Hallazgos relevantes:}
\begin{itemize}
    \item \textbf{Redundancia estructural fuerte:} \texttt{Beat} y
    \texttt{District} tienen una correlación de $r \approx 1.00$, lo cual es
    esperado porque \texttt{Beat} es una subdivisión anidada dentro de
    \texttt{District}: ambas variables aportan prácticamente la misma
    información a distinta granularidad.
    \item \textbf{Relación moderada esperada:} \texttt{Ward} y
    \texttt{Community Area} guardan una correlación moderada
    ($|r|\approx 0.5$--$0.66$) con \texttt{District}/\texttt{Beat}, dado que
    todas son formas alternativas de particionar geográficamente la ciudad.
    \item \textbf{Relación inesperada (casi nula):} \texttt{Hour} es
    prácticamente independiente ($r\approx 0.00$) de todas las variables
    espaciales, lo que sugiere que el patrón horario del crimen es similar
    en toda la ciudad y no varía sustancialmente de un distrito a otro; esto
    debe validarse en la etapa de Visual Analytics, ya que el coeficiente de
    Pearson no captura bien relaciones no lineales o de variables
    cíclicas como la hora.
    \item Las variables \texttt{Arrest} y \texttt{Domestic} muestran
    correlaciones lineales débiles con el resto ($|r|\le 0.10$); esto no
    implica ausencia de relación, sino que dichas relaciones son mejor
    capturadas por pruebas para variables categóricas (ver $\chi^2$ en
    Sección \ref{sec:hipotesis}) que por un coeficiente de correlación
    lineal.
\end{itemize}

% ============================================================================
\section{Problemas identificados y transformaciones (data wrangling)}

\subsection{Problemas detectados en los datos}

\begin{enumerate}
    \item \textbf{Valores nulos en coordenadas geográficas}
    (\texttt{X Coordinate}, \texttt{Y Coordinate}, \texttt{Latitude},
    \texttt{Longitude}, \texttt{Location}): 6,102 registros (3.54\%) sin
    geocodificación.
    \item \textbf{Valores nulos en \texttt{Location Description}}: 515
    registros (0.30\%).
    \item \textbf{Valor nulo en \texttt{Ward}}: 1 registro (0.0006\%).
    \item \textbf{Inconsistencia de dominio en \texttt{District}}: 2
    registros con \texttt{District = 31}, un código administrativo que no
    corresponde a ninguno de los 22 distritos reales de patrullaje de
    Chicago.
    \item \textbf{Inconsistencia de dominio en \texttt{Community Area}}: 1
    registro con valor 0, fuera del rango válido de la ciudad (1--77).
    \item \textbf{Artefacto de captura horaria:} se detectó una
    concentración anómala de reportes exactamente a las 12:00:00~PM: el
    43.5\% de los 10,427 registros de la hora 12 tienen minuto y segundo
    exactamente en cero, frente a solo 31.6\% en la hora de control (13:00).
    Esto sugiere que "mediodía en punto" se usa como marca de tiempo por
    defecto cuando la hora exacta del incidente es desconocida, lo cual
    puede sesgar cualquier análisis fino de patrones horarios si no se
    documenta.
    \item \textbf{Redundancia de variables:} \texttt{Beat} y
    \texttt{District} están altamente correlacionados ($r\approx 1.00$)
    porque \texttt{Beat} es una subdivisión anidada de \texttt{District}.
    \item \textbf{Formato de fecha como texto:} el campo \texttt{Date}
    requiere parseo explícito antes de poder derivar variables temporales.
    \item \textbf{Codificación numérica de variables categóricas:}
    \texttt{District}, \texttt{Ward}, \texttt{Beat} y
    \texttt{Community Area} son códigos categóricos representados como
    números, lo que puede inducir a tratarlos erróneamente como variables
    continuas.
    \item \textbf{Alta cardinalidad y cola larga en variables categóricas:}
    \texttt{Location Description} (123 categorías) e \texttt{IUCR} (318
    categorías) tienen muchas categorías con muy baja frecuencia.
    \item \textbf{Granularidad no uniforme en \texttt{Case Number}:} 27
    números de caso están asociados a más de un registro (2 a 4 filas),
    correspondiendo en el 100\% de los casos a homicidios con múltiples
    víctimas, lo cual debe tenerse en cuenta al contar "incidentes" versus
    "registros".
    \item \textbf{Cobertura temporal parcial:} el extracto usado cubre del
    1 de enero al 25 de agosto de 2017 (no el año calendario completo), por
    lo que las comparaciones de totales brutos por mes deben normalizarse
    por número de días para no favorecer artificialmente a los meses con
    más días (ver H2 en la Sección \ref{sec:hipotesis}).
\end{enumerate}

\subsection{Transformaciones realizadas}

\begin{enumerate}
    \item \textbf{Nulos en coordenadas geográficas:} no se imputaron (crear
    una ubicación ficticia distorsionaría el análisis espacial); se
    excluyen explícitamente solo en los análisis que requieren coordenadas
    (mapa de dispersión, Figura \ref{fig:mapa}), documentando el 3.54\%
    excluido, mientras se conservan para el resto de los análisis.
    \item \textbf{Nulos en \texttt{Location Description}:} se excluyen
    únicamente del análisis de la Hipótesis 4, sin afectar el resto del
    dataset.
    \item \textbf{Registro con \texttt{Ward} nulo:} al tratarse de un único
    registro (0.0006\% del total), se descarta solo en agregaciones por
    \texttt{Ward}, sin impacto material en los resultados.
    \item \textbf{\texttt{District = 31}:} se excluye de los análisis
    agregados por distrito (2 de 172,327 registros) para no introducir una
    categoría espuria en los gráficos espaciales, documentando su existencia
    en vez de eliminarlo del dataset completo.
    \item \textbf{\texttt{Community Area = 0}:} se documenta como "no
    geolocalizado a nivel de área comunitaria" y se excluye únicamente de
    los análisis agregados por dicha variable.
    \item \textbf{Parseo de fecha:} conversión de \texttt{Date} (texto) a
    tipo \texttt{datetime} de pandas; se derivan las nuevas variables
    \texttt{Hour}, \texttt{DayOfWeek} y \texttt{Month} (creación de
    variables / \textit{feature engineering}).
    \item \textbf{Recodificación booleano $\rightarrow$ numérico:}
    \texttt{Arrest} y \texttt{Domestic} se mapean también a 0/1 para poder
    incluirlas en la matriz de correlación y en las pruebas $\chi^2$.
    \item \textbf{Tipificación explícita:} \texttt{District}, \texttt{Ward},
    \texttt{Beat} y \texttt{Community Area} se documentan y tratan
    explícitamente como variables categóricas nominales codificadas,
    evitando interpretarlas como magnitudes continuas pese a que pandas las
    infiere como enteros.
    \item \textbf{Verificación de duplicados:} se comprobó la unicidad de
    \texttt{ID} y de filas completas (0 duplicados encontrados); no se
    requirió deduplicación general. Se documentó por separado el caso de
    \texttt{Case Number} repetido en homicidios (ver arriba), que
    \emph{no} es un error de calidad sino una característica real del
    dominio.
    \item \textbf{Reducción de cardinalidad para visualización:} en los
    gráficos de barras y en el violin plot se seleccionan los "top-N"
    valores más frecuentes (5, 6, 8 o 10, según el gráfico) de las
    variables de alta cardinalidad (\texttt{Primary Type},
    \texttt{Location Description}) para mantener las visualizaciones
    legibles.
    \item \textbf{Normalización temporal para comparación justa entre
    meses:} dado que agosto de 2017 solo tiene 25 días en este extracto
    (frente a 28--31 días de los demás meses), se calculó adicionalmente la
    \textbf{tasa diaria promedio} de delitos por mes (total de delitos /
    número de días con datos) para evitar conclusiones erróneas basadas
    únicamente en el conteo bruto.
\end{enumerate}

% ============================================================================
\section{Exploración y respuesta a las hipótesis}
\label{sec:hipotesis}

La Figura \ref{fig:dashboard} presenta un dashboard resumen con cuatro vistas
generales del dataset, que contextualiza las respuestas detalladas a cada
hipótesis presentadas a continuación.

\begin{figure}[H]
\centering
\includegraphics[width=0.85\textwidth]{../outputs/figures/fig10_dashboard.png}
\caption{Dashboard resumen: delitos por mes, por hora del día, por día de la
semana y los 5 distritos con más delitos. Chicago, 2017.}
\label{fig:dashboard}
\end{figure}

\subsection*{Hipótesis 1: concentración espacial (hotspots)}

\begin{figure}[H]
\centering
\includegraphics[width=0.8\textwidth]{../outputs/figures/fig3_top_crime_types.png}
\caption{Los 10 tipos de delito más frecuentes en Chicago, 2017.}
\label{fig:top_tipos}
\end{figure}

\begin{figure}[H]
\centering
\includegraphics[width=0.72\textwidth]{../outputs/figures/fig6_district_scatter_map.png}
\caption{Distribución geográfica de una muestra aleatoria de 20,000 delitos,
coloreados por distrito policial. Se aprecia claramente la forma de la
ciudad de Chicago y la delimitación de los 22 distritos de patrullaje.}
\label{fig:mapa}
\end{figure}

\textbf{Interpretación.} THEFT (40,742) y BATTERY (32,394) son, por amplio
margen, los tipos de delito más frecuentes (Figura \ref{fig:top_tipos}). En
cuanto a distribución espacial, los 5 distritos con más delitos --11
(11,778), 6 (10,754), 8 (10,698), 1 (9,748) y 4 (9,437)-- concentran el
\textbf{30.4\%} del total de incidentes (excluyendo el código administrativo
atípico \texttt{District 31}), mientras que el distrito con menos delitos
(20) registra solo 2,804. La Figura \ref{fig:mapa} confirma visualmente que
el crimen no se distribuye uniformemente en el espacio urbano.

\textbf{Conclusión intermedia.} Se \textbf{confirma} la Hipótesis 1: existe
una concentración espacial desigual del crimen (aproximadamente un tercio de
los delitos ocurre en solo 5 de los 23 distritos), consistente con la
existencia de hotspots espaciales, tal como plantea CrimAnalyzer para São
Paulo. La concentración es notoria pero no extrema, lo que sugiere que un
enfoque como el de CrimAnalyzer (basado en NMF, capaz de detectar también
hotspots "menores pero recurrentes") aportaría más valor que un simple
ranking por volumen.

\subsection*{Hipótesis 2: patrones temporales}

\begin{figure}[H]
\centering
\includegraphics[width=0.75\textwidth]{../outputs/figures/fig1_hour_histogram.png}
\caption{Distribución de delitos según la hora del día. Chicago, 2017.}
\label{fig:hist_hora}
\end{figure}

\begin{figure}[H]
\centering
\includegraphics[width=0.68\textwidth]{../outputs/figures/fig5_hour_weekday_heatmap.png}
\caption{Delitos por hora del día y día de la semana. Chicago, 2017.}
\label{fig:heatmap_hora_dia}
\end{figure}

\begin{figure}[H]
\centering
\includegraphics[width=0.85\textwidth]{../outputs/figures/fig8_hour_violin_by_type.png}
\caption{Distribución horaria según los 5 tipos de delito más frecuentes.}
\label{fig:violin}
\end{figure}

\textbf{Interpretación.} La hora con más delitos es las 12:00 (10,427
casos) y la de menos es las 5:00 de la madrugada (2,511 casos) --Figura
\ref{fig:hist_hora}--. Sin embargo, como se documentó en la Sección 7, el
pico de las 12:00 está \textbf{inflado por un artefacto de captura}
(uso de "mediodía en punto" como marca por defecto), por lo que el patrón
horario real probablemente crece de forma más gradual desde la mañana hasta
la noche, como sugiere el heatmap de la Figura \ref{fig:heatmap_hora_dia}:
la franja de 12:00 a 20:00 concentra visiblemente más actividad en todos los
días de la semana, sin diferencias notorias entre día laborable y fin de
semana. Por día de la semana, la variación es baja (entre 24,204 y 25,813
delitos), con un ligero máximo el viernes. Al comparar meses mediante la
tasa diaria promedio (para corregir por la duración desigual de los meses,
dado que agosto solo tiene 25 días en este extracto), julio resulta el mes
con mayor intensidad (781.7 delitos/día) y marzo el de menor intensidad
(654.1 delitos/día) --una conclusión distinta a la que se obtendría
comparando solo totales brutos, donde febrero aparentaría ser el más bajo
simplemente por tener menos días--. La Figura \ref{fig:violin} muestra
además que distintos tipos de delito tienen "firmas" horarias distintas
(por ejemplo, robos y hurtos muestran mayor densidad vespertina que otros
tipos).

\textbf{Conclusión intermedia.} Se \textbf{confirma} la Hipótesis 2: existen
patrones temporales claros y estadísticamente relevantes tanto a nivel de
hora del día como de mes, aunque el día de la semana influye poco. Estos
patrones horarios varían además según el tipo de delito, lo que refuerza la
necesidad de vistas temporales interactivas (no un solo agregado) en la
etapa de Visual Analytics.

\subsection*{Hipótesis 3: tasa de arresto según tipo de delito}

\begin{figure}[H]
\centering
\includegraphics[width=0.75\textwidth]{../outputs/figures/fig7_arrest_rate_by_type.png}
\caption{Tasa de arresto según tipo de delito (8 tipos más frecuentes).}
\label{fig:arresto}
\end{figure}

\begin{figure}[H]
\centering
\includegraphics[width=0.65\textwidth]{../outputs/figures/fig2_hour_by_arrest_boxplot.png}
\caption{Hora del delito según si hubo arresto.}
\label{fig:boxplot}
\end{figure}

\textbf{Interpretación.} La tasa de arresto global es 19.24\%, pero varía
fuertemente según el tipo de delito (Figura \ref{fig:arresto}): "OTHER
OFFENSE" tiene la tasa más alta (22.2\%), seguida de BATTERY (19.8\%) y
ASSAULT (17.8\%); en el otro extremo, BURGLARY y DECEPTIVE PRACTICE
comparten la tasa más baja (3.7\%). La prueba $\chi^2$ de independencia
entre \texttt{Primary Type} y \texttt{Arrest} (8 categorías más frecuentes)
da $\chi^2 = 5476.7$, $\text{gl} = 7$, $p < 0.001$, rechazando con amplio
margen la hipótesis nula de independencia. El boxplot de la Figura
\ref{fig:boxplot} no muestra una diferencia horaria marcada entre incidentes
con y sin arresto.

\textbf{Conclusión intermedia.} Se \textbf{confirma} la Hipótesis 3: el tipo
de delito es un fuerte predictor de si el incidente termina en arresto. La
categoría "OTHER OFFENSE" -heterogénea por definición- concentra la mayor
tasa, lo que sugiere revisar en el sistema de Visual Analytics si conviene
desagregarla usando el campo \texttt{Description} o \texttt{IUCR} para no
perder interpretabilidad.

\subsection*{Hipótesis 4: violencia doméstica según tipo de lugar}

\begin{figure}[H]
\centering
\includegraphics[width=0.75\textwidth]{../outputs/figures/fig9_domestic_rate_by_location.png}
\caption{Proporción de incidentes domésticos según tipo de lugar (6 lugares
más frecuentes).}
\label{fig:domestico}
\end{figure}

\textbf{Interpretación.} La tasa global de incidentes marcados como
\texttt{Domestic} es 16.31\%, pero se dispara a 43.1\% en
\texttt{APARTMENT} y 33.8\% en \texttt{RESIDENCE}, frente a solo 8.2\% en
\texttt{STREET} y 5.8\% en \texttt{PARKING LOT/GARAGE} (Figura
\ref{fig:domestico}). La prueba $\chi^2$ de independencia entre
\texttt{Location Description} y \texttt{Domestic} (6 categorías más
frecuentes) da $\chi^2 = 14917.8$, $\text{gl} = 5$, $p < 0.001$, un efecto
aún más marcado que en la Hipótesis 3.

\textbf{Conclusión intermedia.} Se \textbf{confirma} la Hipótesis 4: la
violencia doméstica está fuertemente asociada al hogar como lugar de
ocurrencia (apartamento/residencia), en línea con el conocimiento de
dominio en criminología. Este resultado también funciona como una
verificación de consistencia (\textit{sanity check}) que respalda la
validez general del dataset.

% ============================================================================
\section{Conclusiones}

\textbf{Resumen del conocimiento obtenido.} El dataset de crímenes de
Chicago (extracto ene--ago 2017, 172,327 registros) exhibe: (i) una
distribución de tipos de delito muy sesgada, dominada por THEFT y BATTERY;
(ii) una concentración espacial moderada del crimen en unos pocos distritos
(30.4\% en los 5 principales); (iii) un ritmo temporal diario marcado, con
mayor actividad entre el mediodía y la noche y mínima actividad de
madrugada, y una intensidad mensual algo mayor en los meses de verano
boreal (julio); (iv) diferencias grandes y estadísticamente significativas
en la tasa de arresto según el tipo de delito, y (v) una fuerte asociación,
también significativa, entre violencia doméstica y el hogar como lugar del
incidente.

\textbf{Respuesta final a las hipótesis.}
\begin{itemize}
    \item \textbf{H1 (espacial): confirmada.} Los 5 distritos con más
    delitos concentran 30.4\% del total; existe evidencia de hotspots
    espaciales.
    \item \textbf{H2 (temporal): confirmada.} Hay patrones horarios y
    mensuales claros (aunque el patrón horario en punta requiere lectura
    cuidadosa por el artefacto de captura a las 12:00), con muy poca
    variación por día de la semana.
    \item \textbf{H3 (esclarecimiento): confirmada.} La tasa de arresto
    depende fuertemente del tipo de delito ($\chi^2=5476.7$, $p<0.001$).
    \item \textbf{H4 (violencia doméstica y lugar): confirmada.} La
    proporción de casos domésticos depende fuertemente del tipo de lugar
    ($\chi^2=14917.8$, $p<0.001$).
\end{itemize}

\textbf{Principales sorpresas y desafíos.} El hallazgo más relevante desde
el punto de vista de calidad de datos fue el \textbf{artefacto de captura a
las 12:00:00~PM}, que de no documentarse podría llevar a sobreestimar la
criminalidad al mediodía. También resultó revelador que el extracto
\emph{no} cubre el año calendario completo (termina el 25 de agosto), lo
que obliga a normalizar por días antes de comparar meses; y que la
aparente "duplicidad" de \texttt{Case Number} en 27 casos no es un error,
sino el reflejo correcto de que los homicidios se registran por víctima.
Finalmente, la fuerte redundancia entre \texttt{Beat} y \texttt{District}
($r\approx 1.00$) es un recordatorio de que no toda variable adicional
aporta información nueva.

\textbf{Puente hacia Visual Analytics.} Estos hallazgos definen las
prioridades de la siguiente etapa del curso, el diseño del sistema de Visual
Analytics inspirado en CrimAnalyzer \cite{crimanalyzer}:

\begin{enumerate}
    \item \textbf{Vistas espacio-temporales enlazadas:} un mapa interactivo
    (a nivel de distrito o de coordenadas) enlazado con una vista temporal
    (hora $\times$ día, similar a la Figura \ref{fig:heatmap_hora_dia}),
    replicando la filosofía central de CrimAnalyzer de exploración conjunta
    espacio + tiempo.
    \item \textbf{Detección de hotspots más allá del volumen bruto:} dado
    que la concentración espacial es notoria pero no extrema, conviene
    implementar un mecanismo de detección de hotspots (p. ej. NMF, como en
    CrimAnalyzer) que identifique zonas de interés que un ranking simple por
    volumen pasaría por alto.
    \item \textbf{Filtros interactivos por tipo de delito, arresto y lugar:}
    justificados por las fuertes diferencias encontradas en las Hipótesis
    3 y 4.
    \item \textbf{Manejo explícito en la interfaz de la incertidumbre de los
    datos:} el sistema debe señalar visualmente los registros sin
    geocodificar (3.54\%) y advertir sobre el artefacto de las 12:00~PM, en
    vez de ocultarlos silenciosamente.
    \item \textbf{Evitar redundancia visual:} no usar \texttt{Beat} y 
    \texttt{District} simultáneamente como canales visuales independientes,
    dada su altísima correlación.
\end{enumerate}

% ============================================================================
\bibliographystyle{plainnat}
\bibliography{refs}

% ============================================================================
\section{Anexos}

\begin{itemize}
    \item \textbf{Notebook completo y reproducible:}
    \texttt{analisis\_chicago\_2017.ipynb} (adjunto a esta entrega), que
    contiene la carga de datos, el wrangling, el cálculo de todos los
    estadísticos citados en este informe y el código de generación de las
    10 figuras utilizadas.
    \item \textbf{Datos:} archivo \texttt{crimes\_2017.csv} (extracto
    oficial enero--agosto 2017 usado en este informe) y enlace al dataset
    completo en el portal oficial (Sección 4).
    \item \textbf{Figuras:} carpeta \texttt{../outputs/figures/}, generada
    automáticamente al ejecutar el notebook.
\end{itemize}

\end{document}